

\begin{derivation}[从\eqnrefs{eq:once-subtracted-single-cut-dp9}到\eqnrefs{eq:spacelike-spectral-positivity-dp68}]
这一结论只讨论谱密度 $\rho(s)\equiv\operatorname{Im}F(s+\ii0^+)\ge0$ 的投影。令 $z=-Q^2$、$Q^2>0$，正文\eqnrefs{eq:once-subtracted-single-cut-dp9} 变为
\begin{equation*}
F(-Q^2)
=-\frac{Q^2}{\pi}
\int_{s_0}^{\infty}\dd s'\,
\frac{\rho(s')}{s'(s'+Q^2)}.
\end{equation*}
利用恒等式
\begin{equation*}
\frac{Q^2}{s'(s'+Q^2)}
=\frac1{s'}-\frac1{s'+Q^2},
\end{equation*}
可以把所有 $Q^2$ 依赖集中到第二项：
\begin{equation*}
F(-Q^2)
=-\frac1\pi\int_{s_0}^{\infty}\frac{\rho(s')}{s'}\dd s'
+\frac1\pi\int_{s_0}^{\infty}\frac{\rho(s')}{s'+Q^2}\dd s'.
\end{equation*}
若相应谱矩收敛，可以逐次把 $Q^2$ 导数移入积分号。一次和二次导数分别为
\begin{align*}
\frac{\dd}{\dd Q^2}F(-Q^2)
&=-\frac1\pi\int_{s_0}^{\infty}\dd s'\,
\frac{\rho(s')}{(s'+Q^2)^2}\le0,\\
\frac{\dd^2}{\dd(Q^2)^2}F(-Q^2)
&=\frac{2}{\pi}\int_{s_0}^{\infty}\dd s'\,
\frac{\rho(s')}{(s'+Q^2)^3}\ge0.
\end{align*}
继续求导时，每次只会再产生一个负号和一个正的整数系数，因此归纳得到正文\eqnrefs{eq:spacelike-spectral-positivity-dp68}。它给出了一个很实用的自洽检查：若显式圈积分声称对应同一非负谱投影，却违反这一导数符号层级，就说明符号、减法或解析延拓中至少有一处不一致。
\end{derivation}
